\section{Claims Mapping}
\label{app:claims}

This appendix pins every body claim to one of three audit buckets. \texttt{supported} claims have direct artifact evidence (\texttt{allclose} or exact-token-match). \texttt{proxy\_supported} claims have proxy-attacker or cost-model evidence only and are worded ``in our tested configurations.'' \texttt{unsupported} claims must \emph{never} appear as positive claims in the body.

\subsection*{Supported claims (S1--S8)}
\begin{description}[leftmargin=*, style=unboxed]
  \item[S1] GPT-2 model-level masked execution reproduces the plain reference output token-for-token in our tested configurations.
  \item[S2] Modern decoder-only model-level masked execution reproduces the plain reference output token-for-token in our tested configurations.
  \item[S3] Compatible mask families preserve correctness across the tested architectures.
  \item[S4] KV cache append invariant holds for the tested decoder-only generation paths.
  \item[S5] Masked LoRA forward/backward/rank-padded primitives match the plain rank-$\rtrue$ reference to float64 precision on synthetic single-linear tests.
  \item[S6] Synthetic multi-layer rank-padded LoRA training matches the plain reference on every per-module slice.
  \item[S7] Every tested stronger-dummy strategy preserves the rank-padded LoRA invariant $\Apad\Bpad = \A\B$ (exactly or via a tracked trusted-side correction) and matches plain training to float64 precision.
  \item[S8] Under the cost-model proxy, equalizing per-step latency to the upper bucket reduces the worst-case timing classifier accuracy to near random chance.
\end{description}

\subsection*{Proxy-supported claims (P1--P5)}
\begin{description}[leftmargin=*, style=unboxed]
  \item[P1] Under our adaptive proxy attackers (ridge, MLP, signature- and Sinkhorn-based permutation recovery, linkability), the full mitigation bundle keeps the worst-case attacker close to random chance in our tested configurations.
  \item[P2] Under our cost-model timing proxy, the \texttt{proxy\_equalized} mode reduces classifier accuracy to near random chance.
  \item[P3] Under our LoRA adapter and gradient leakage proxy attackers, fresh masks $+$ pad bring the linkability AUC close to $0.5$ in our tests.
  \item[P4] Rank padding hides $\rtrue$ from tensor shape; under our spectral-cliff, energy, elbow, and ensemble proxy detectors, stronger dummy strategies fail to recover $\rtrue$ reliably in our tests (\texttt{needs\_more\_evaluation}).
  \item[P5] Under our cross-layer linkability proxy, fresh masks per module keep the AUC near $0.5$ across the tested multi-layer configuration.
\end{description}

\subsection*{Unsupported claims (U1--U8) --- must NOT appear as positive claims in the body}
\begin{description}[leftmargin=*, style=unboxed]
  \item[U1] Formal/cryptographic/semantic security of the masked path. Safe wording: \emph{We make no formal/cryptographic/semantic security claims.}
  \item[U2] Real TEE wall-time. Safe wording: \emph{We report measured local-emulation latencies; we do not claim real TEE wall-time.}
  \item[U3] Hardware side-channel security (cache, power, EM).
  \item[U4] Full Qwen/TinyLlama/LLaMA LoRA fine-tuning.
  \item[U5] PEFT/DeepSpeed/vLLM/FlashAttention compatibility.
  \item[U6] $\rpad$ is hidden from the GPU.
  \item[U7] Loss/optimizer is fully outsourced to the GPU.
  \item[U8] Protection against a compromised TEE.
\end{description}

\subsection*{Body $\to$ claim mapping}
\Cref{sec:introduction} (Contributions) maps C1 to S1, S2, S4; C2 to S3; C3 to S5, S6, S7; C4 cites the audit and summary directly. \Cref{sec:threat} lists U1--U8 verbatim under out-of-scope. \Cref{sec:design:rank-pad,sec:design:lora-backward} explicitly call out U6 (padded rank visible) and U7 (loss/optimizer trusted). \Cref{sec:correctness} carries only S1--S8 under theorem labels. \Cref{sec:security} carries only P1--P5 in positive statements; every paragraph is hedged with ``in our tested configurations.'' \Cref{sec:evaluation} tags each RQ with the claims it instantiates. \Cref{sec:limitations} expands U1--U8 into items 1--15. \Cref{sec:related} uses ``proxy-evaluated, not formal'' throughout. \Cref{sec:conclusion} closes with the same U1--U8 disclaimer block.

\paragraph{Wording check.}
If a candidate sentence contains any of: \texttt{secure}, \texttt{provably}, \texttt{cryptographic}, \texttt{semantic security}, \texttt{prevents all leakage}, \texttt{guarantees}, \texttt{real TEE wall-time}, \texttt{hides padded rank}, \texttt{Qwen fine-tune}, \texttt{LLaMA fine-tune}, \texttt{PEFT}, \texttt{DeepSpeed}, \texttt{vLLM}, \texttt{FlashAttention}, or \texttt{fully outsourced loss/optimizer}, that sentence requires a corresponding \texttt{unsupported} audit row; re-read U1--U8 and either re-word or move to Limitations. The complete audit lives in \texttt{paper\_results/markdown/paper\_claims\_audit.md} and is reproduced verbatim in \Cref{tab:paper-claims-audit}.

% Auto-generated by paper_claims_audit
\begin{table}[h]
\centering
\small
\caption{Paper claims audit: supported, proxy-supported, unsupported.}
\label{tab:paper_claims_audit}
\begin{tabular}{p{1.6cm}p{6cm}p{6cm}}
\toprule
\textbf{Status} & \textbf{Claim} & \textbf{Paper-safe wording} \\
\midrule
supported & Masked execution preserves output equality vs the plain reference for the tested GPT-2 model-level wrapper (full forward + prefill + decode\_step + greedy generation). & We empirically verify that the GPT-2 model-level masked wrapper reproduces the plain reference output token-for-token in our tested configurations. \\
supported & Masked execution preserves output equality vs the plain reference for the tested modern-decoder model-level wrapper (RMSNorm + SwiGLU + RoPE + GQA, prefill + decode\_step + greedy generation). & We empirically verify that the modern decoder-only model-level masked wrapper reproduces the plain reference output token-for-token in our tested configurations. \\
supported & Compatible nonlinear islands preserve output correctness across the tested architectures (decoder-only / encoder-only / encoder-decoder, plus modern decoder). & Compatible mask families (mean-preserving orthogonal for LayerNorm, channel permutation for GELU/ReLU, paired permutation for SwiGLU) preserve output correctness in our tests; we do not claim they extend beyond the tested mask families. \\
supported & KV cache append invariant holds for the tested decoder-only generation paths. & We verify the KV cache append invariant matches the plain reference in our tested decoder-only generation paths. \\
supported & LoRA masked forward, masked backward, and rank-padded forward/backward reproduce the plain rank-r reference to float64 precision in synthetic single-linear tests. & We empirically verify that the masked LoRA forward / backward / rank-padded primitives match the plain rank-r reference to float64 precision on synthetic single-linear tests. \\
supported & Multi-layer rank-padded LoRA training reproduces the plain reference loss / gradient / SGD / AdamW update across every per-module slice in synthetic multi-layer tests. & We empirically verify that a synthetic multi-layer rank-padded LoRA training step matches the plain reference on every per-module slice. \\
supported & All Stage 7.4 stronger dummy distributions preserve A\_pad B\_pad = A\_real B\_real exactly (cancellation strategies) or with a tracked trusted-side correction (noise\_injected\_cancellation\_dummy), and all five preserve forward / backward / SGD  & Every tested stronger-dummy strategy preserves the rank-padded LoRA invariant A\_pad B\_pad = A\_real B\_real (either exactly or via a tracked trusted-side correction) and matches plain training to float64 precision in our tests. \\
supported & Constant-time training mode (proxy\_equalized) reduces the cost-model timing classifier accuracy to near random chance in our tests. & Under the cost-model proxy, equalizing per-step latency to the upper bucket reduces the worst-case timing classifier accuracy to near random chance in our tests. \\
proxy\_supported & Activation-recovery / linkability adaptive proxy attackers are bounded under the full mitigation bundle (fresh permutation + dense sandwich + boundary pad) in the tested configurations. & Under our adaptive proxy attackers (ridge, small MLP, signature / Sinkhorn permutation recovery, linkability), the full mitigation bundle keeps the worst-case attacker close to random chance in our tested configurations. \\
proxy\_supported & Cost-model timing side-channel proxy is bounded under proxy\_equalized constant-time decode / training modes in our tests. & Under our cost-model timing proxy, equalizing latency to the upper bucket reduces classifier accuracy near random chance. \\
proxy\_supported & LoRA adapter extraction / gradient leakage proxy attackers are bounded under fresh masks + pad in the tested configurations. & Under our LoRA adapter / gradient leakage proxy attackers, fresh masks + pad bring the linkability AUC close to 0.5 (random chance) in our tests. \\
proxy\_supported & Rank padding hides true\_rank from the GPU-visible tensor shape, and stronger dummy distributions reduce the proxy spectral-cliff inference accuracy of true\_rank. & Rank padding hides true\_rank from tensor shape; under our spectral-cliff / energy / elbow / ensemble proxy detectors, stronger dummy strategies fail to recover true\_rank reliably in our tests (`needs\_more\_evaluation`). \\
proxy\_supported & Cross-layer adapter linkage is low under fresh masks + paired cancellation across the tested multi-layer configuration. & Under our cross-layer linkability proxy, fresh masks per module keep the AUC near 0.5 across the tested multi-layer configuration. \\
unsupported & Formal / cryptographic / semantic security of the masked path. & We make no formal / cryptographic / semantic security claims. \\
unsupported & Real TEE wall-time / hardware-isolation evaluation. & We report measured local-emulation latencies; we do not claim real TEE wall-time. \\
unsupported & Hardware side-channel security (cache / power / EM). & We do not evaluate hardware side-channels (cache / power / EM). \\
unsupported & Full Qwen / TinyLlama / LLaMA LoRA fine-tuning. & Our LoRA experiments use synthetic single-linear and multi-layer tiles; we do not run full Qwen / TinyLlama / LLaMA LoRA fine-tuning. \\
unsupported & PEFT / DeepSpeed / vLLM / FlashAttention compatibility. & We do not integrate PEFT / DeepSpeed / vLLM / FlashAttention; our LoRA primitives are stand-alone functional API. \\
unsupported & padded\_rank is hidden from the GPU. & Stage 7.2 / 7.3 / 7.4 hide true\_rank from tensor shape; padded\_rank itself remains visible to the GPU. \\
unsupported & Loss / optimizer is fully outsourced to the GPU. & Loss + optimizer remain trusted-side; only forward / backward matmuls cross the boundary. \\
unsupported & Protection against a compromised TEE. & We assume the TEE is honest; a compromised TEE breaks every guarantee we evaluate. \\
\bottomrule
\end{tabular}
\end{table}

